% ID: SEL-D0028-OB-E03
% ESTADO: validada. S029, cotejada con examen y criterios.
\textbf{a)} Escribimos $N$, $B$ y $R$ para negra, blanca y roja.
Importa el orden y sólo hay una bola roja. Por tanto
$\Omega=\{NN,NB,NR,BN,BB,BR,RN,RB\}$.

\textbf{b)} Si distinguimos las bolas físicas, hay $6\cdot5=30$
pares ordenados equiprobables. Para cada par de colores,
multiplicamos las posibilidades de la primera extracción por
las que quedan para la segunda.
\[
\begin{array}{c|rrrrrrrr}
\text{Resultado}&NN&NB&NR&BN&BB&BR&RN&RB\\\hline
\text{Casos}&6&6&3&6&2&2&3&2\\
\text{Probabilidad}&6/30&6/30&3/30&6/30&2/30&2/30&3/30&2/30
\end{array}
\]
La suma de los casos es 30 y la de las probabilidades es 1.
